\documentclass[12pt]{article}
\usepackage[utf8]{inputenc}
\usepackage[T1]{fontenc}
\usepackage[margin=1in]{geometry}
\usepackage{setspace}
\usepackage{xcolor}
\usepackage[colorlinks=true, urlcolor=blue!60!black, linkcolor=blue!60!black]{hyperref}
\setlength{\parindent}{2em}
\setlength{\parskip}{1.25em}
\renewcommand{\baselinestretch}{1.0}
\begin{document}
\begin{center}
{\LARGE \textit{Agency for International Development v.\ Alliance for Open Society International, Inc.}, 570 U.S.\ 205 (2013)\par}
\vspace{0.8em}
{\large Case Brief --- Four-Page Edition\par}
\end{center}
\vspace{1em}

\section{Memory Jogger}

Under the Spending Clause and the unconstitutional conditions doctrine, Congress may specify the activities it wishes to subsidize, but it cannot compel domestic grant recipients to adopt and espouse the Government's viewpoint as their own belief outside the contours of the funded program (\href{https://supreme.justia.com/cases/federal/us/570/205/}{570 U.S. at 214--15, 221}).

\section{Detailed Case Facts}

The Leadership Act of 2003 authorized billions of dollars for nongovernmental organizations fighting HIV/AIDS abroad and declared a national policy to eradicate prostitution. It attached two conditions: an unchallenged ban on using funds to promote or advocate the legalization or practice of prostitution (\S~7631(e)), and a Policy Requirement barring funds to any group that does not have a policy explicitly opposing prostitution and sex trafficking (\S~7631(f)). HHS and USAID required every recipient to state that opposition in its award documents. Respondents, domestic NGOs such as AOSI and Pathfinder International, did not promote prostitution but wished to remain neutral: condemning it risked alienating host governments, losing the trust of sex workers needed for HIV prevention, and censoring their privately funded advocacy.

\section{Procedural History}

AOSI and Pathfinder sued in the Southern District of New York in 2005 and won a preliminary injunction. While the appeal was pending, the agencies issued affiliate guidelines letting recipients work with separate affiliates that spoke differently. The district court reissued the injunction, and the Second Circuit affirmed (2--1). Because the D.C. Circuit had upheld \S~7631(f) in \textit{DKT International}, the Supreme Court granted certiorari to resolve the split.

\section{Judicial Votes}

Affirmed, \textbf{6--2} (decided June 20, 2013).
\begin{itemize}
\item \textbf{Majority:} Chief Justice Roberts, joined by Justices Kennedy, Ginsburg, Breyer, Alito, and Sotomayor.
\item \textbf{Dissent:} Justice Scalia, joined by Justice Thomas.
\item \textbf{Recusal:} Justice Kagan took no part.
\end{itemize}

\section{Holding}

The Policy Requirement violates the First Amendment as applied to domestic organizations: Congress may not condition funding on a recipient's affirmative adoption of a government viewpoint that, by its nature, cannot be confined within the funded program. The Court distinguished permissible conditions, which \emph{define the limits of the government spending program}, from impermissible ones, which \emph{leverage funding to regulate speech outside the contours of the program itself}. The Court did not decide whether the requirement could be applied to foreign NGOs.

\section{Analysis of Opinions}

\textbf{Majority (Roberts).} Freedom of speech prohibits the government from telling people what they must say; as a direct regulation, the Policy Requirement would plainly be invalid. Synthesizing precedent, Roberts drew the inside/outside line: \textit{Regan v. Taxation With Representation} upheld a lobbying limit because a dual affiliate could lobby with private funds; \textit{FCC v. League of Women Voters} struck down a total editorializing ban that reached privately funded speech; and \textit{Rust v. Sullivan} upheld abortion-counseling limits confined to the Title X \emph{project}, not the grantee. Because \S~7631(e) already kept federal money from promoting prostitution, \S~7631(f) necessarily went further, ``defining the recipient.'' Roberts rejected three defenses: the affiliate guidelines, because a separate affiliate cannot express the recipient's own beliefs and a closely identified one forces ``evident hypocrisy''; the money-is-fungible theory, which lacked record support and would overturn \textit{League of Women Voters}; and preventing message confusion, which could not justify forcing grantees to pledge allegiance to the Government's policy, contrary to \textit{Barnette}.

\textbf{Scalia (dissenting).} The Policy Requirement is a commonsense selection criterion letting the Government enlist ideological allies: the United States need not hire Hamas to distribute food aid. Because eradicating prostitution was a core objective of the Act, outside promotion undermined the program itself. Above all, there was no coercion: ``the Government is not forcing anyone to say anything,'' since any NGO may decline the grant. The Oath Clause shows the Framers accepted relevant ideological commitments as a condition of assisting the Government.

\textbf{Subsequent development.} In \textit{AOSI II}, 591 U.S. 430 (2020), the Court held (5--3) that the Policy Requirement may be applied to plaintiffs' separately incorporated foreign affiliates, because foreign entities abroad possess no First Amendment rights.

\section{Examples: Future Applications}

\begin{enumerate}
\item \textbf{Same-Side (Unconstitutional): Opioid Clinics.} Federal overdose-prevention grants are denied to any organization lacking a corporate policy opposing decriminalization of controlled substances. Even with a ban on lobbying with grant money, the creed compels institutional belief outside the program.
\item \textbf{Same-Side (Unconstitutional): University Pledges.} A state conditions research grants on a university's public pledge endorsing the legislature's stance on an international boycott---orthodoxy unrelated to the research subsidized.
\item \textbf{Opposite-Side (Permissible): Clean Energy Curriculum.} A federal youth-corps program requires funded curriculum to teach that solar and wind power reduce carbon emissions, leaving grantees free to publish skeptical articles with private funds. The condition defines the program, as in \textit{Rust}.
\item \textbf{Opposite-Side (Permissible): Foreign NGOs.} An anti-trafficking pledge required of independently incorporated foreign NGOs operating overseas is valid under \textit{AOSI II}.
\item \textbf{Fence-Sitter: Anti-Smoking Grants.} An agency scores applicants on organizational commitment to tobacco cessation and disqualifies an NGO that accepts tobacco-industry sponsorships elsewhere. Is it selecting a credible messenger, or penalizing off-program speech and association?
\end{enumerate}

\section{Critique}

\textbf{Indeterminacy.} Roberts conceded that the program/outside-the-program line is ``hardly clear,'' because Congress can often redefine a program broadly enough to subsume a condition. The Court also never answered Scalia's sharpest challenge: if the Government cannot require applicants to \emph{adopt} a view, may it still \emph{prefer} those who already hold it? If so, the rule elevates form over substance.

\textbf{Competing perspectives.} Scalia and Thomas read ``abridging'' to exclude refusals to subsidize, while libertarian originalists applaud the result for checking the spending power. Progressive critics praise the protection of public-health NGOs but note that \textit{AOSI II} later forced U.S. NGOs, through their foreign affiliates, into the very ``evident hypocrisy'' Roberts condemned, and warn that an expansive compelled-speech doctrine could be turned against nondiscrimination conditions on grants.

\section{Key Quotations}

\begin{quote}
``In the present context, the relevant distinction that has emerged from our cases is between conditions that define the limits of the government spending program---those that specify the activities Congress wants to subsidize---and conditions that seek to leverage funding to regulate speech outside the contours of the program itself.''
\end{quote}
\noindent --- Chief Justice Roberts, for the Court, 570 U.S. at 214--15.

\begin{quote}
``By demanding that funding recipients adopt---as their own---the Government's view on an issue of public concern, the condition by its very nature affects `protected conduct outside the scope of the federally funded program.' \ldots By requiring recipients to profess a specific belief, the Policy Requirement goes beyond defining the limits of the federally funded program to defining the recipient.''
\end{quote}
\noindent --- Chief Justice Roberts, for the Court, 570 U.S. at 218.

\begin{quote}
``Affiliates cannot serve that purpose when the condition is that a funding recipient espouse a specific belief as its own. If the affiliate is distinct from the recipient, the arrangement does not afford a means for the recipient to express its beliefs. If the affiliate is more clearly identified with the recipient, the recipient can express those beliefs only at the price of evident hypocrisy.''
\end{quote}
\noindent --- Chief Justice Roberts, for the Court, 570 U.S. at 219.

\end{document}
